\chapter{Measuring the Theoretical Quantities: What Could and Could Not Be Measured}
\label{ch:instantiation}

\section{Purpose of this chapter}

This chapter tabulates which of the symbols introduced in Part~\ref{part:theory} could
actually be measured, and what values they take where they could. Each chapter of
Part~\ref{part:results} treats a separate domain, so the correspondence with the theoretical
quantities is easily scattered. This chapter supplies the index.

\textbf{Its main purpose is to make explicit what could not be measured.} Fewer than half the
quantities the framework defines could be connected to real data.

The stage reached differs by quantity (Remark~\ref{rem:stage-per-quantity}). The list here is
a record of where between stages 1 and 3 of Section~\ref{sec:five-stages} each one stalled.

\section{The primitive sets made concrete}

\begin{center}
\begin{tabularx}{\textwidth}{@{}llX@{}}
\toprule
Symbol & Domain & The set realized, and its cardinality \\
\midrule
$N$ & family 2 & $\{$issuers, purchasers, merchants$\}$. 2{,}034 issuers (829 third-party, 1{,}205 own-use) \\
$N$ & family 5 & $\{$agencies, job seekers, employers$\}$. 30{,}561 establishments, 13{,}946 with placements \\
$N$ & industry & the population of the corporate statistics: 62 industries $\times$ 16 capital classes \\
$T$ & family 2 & $\{$FY2020--FY2024$\}$, $|T|=5$ \\
$T$ & family 5 & FY2024 alone; the panel covers FY2020--FY2025, $|T|=6$ \\
$T$ & industry & FY2000--FY2024, $|T|=25$ \\
$\mathcal{D}$ & family 2 & four media of prepaid instruments (paper, magnetic, IC, server) \\
$\Omega$ & all domains & \textbf{not observed}. Only the arrival rate is proxied in Chapter~\ref{part:unmanned} \\
$\mathcal{F}_t$ & all domains & \textbf{not observed} \\
\bottomrule
\end{tabularx}
\captionof{table}{How the primitive sets $N$, $T$, $\mathcal{D}$, $\Omega$ and $\mathcal{F}_t$ were realized in each domain.}
\label{tab:basic-sets}
\end{center}

\begin{remark}[Consequence of $\Omega$ being unmeasurable]
The state space itself cannot be observed. The dependence of $\pi$ or $\delta$ on $\omega$
therefore cannot be verified directly, and every empirical result here concerns quantities
marginalized over $\omega$. The one exception is Chapter~\ref{part:unmanned}, where the
arrival rate $\lambda_{\rm novel}$ of unforeseen states is proxied by counts of
vulnerabilities.
\end{remark}

\section{The observed range of the map \texorpdfstring{$\Phi$}{Phi}}

Part~\ref{part:catalog} enumerated 28 types; those that could be connected to real data are
limited to the following (Table~\ref{tab:phi-observed-types}).

\begin{center}
\begin{tabularx}{\textwidth}{@{}llX@{}}
\toprule
Type & Instance & Quantities observed \\
\midrule
2-4 prepaid & prepaid payment instruments & $D$, $P$, $\kappa$, an upper bound on $\dbar-\delta$ \\
5-1 success fee & fee-charging employment placement & $\pi$ per placement, $y$, the extent of $\mathcal{G}$ \\
6-1 transaction fee & payment and distribution platforms & the price of $\phi_{\rm barg}$ \\
6-4 charging for opportunity & marketplace customer acquisition & the price of $T_{\rm acq}$ \\
1-3 one-off arrears & trade credit in the corporate statistics & $\kappa_C>0$ \\
\bottomrule
\end{tabularx}
\captionof{table}{The five of the 28 types that could be connected to real data, and the quantities observed.}
\label{tab:phi-observed-types}
\end{center}

That is \textbf{five of the 28 types}. Family 3 (renting assets), family 4 (complement
recovery) and family 7 (payer separation) were not measured at all.

\section{Cumulants and the credit position}

\subsection{\texorpdfstring{$D$}{D}, \texorpdfstring{$P$}{P} and \texorpdfstring{$\kappa$}{kappa} in family 2}

Substituting FY2024 values into the definitions of Chapter~\ref{sec:phi}, in millions of yen:

\begin{align}
P(\text{FY2024}) &= 28{,}150{,}350 \qquad \text{(issued: settlement from customers)} \\
D(\text{FY2024}) &= 28{,}241{,}983 \qquad \text{(redeemed: delivery by the firm)} \\
-\kappa_C(\text{FY2024}) &= 2{,}895{,}327 \qquad \text{(unused balance: customers credit the firm)}
\end{align}

Against~\eqref{eq:kappa}, $\kappa_C = -2.90$ trillion yen. Being negative, customers extend
credit to the issuer. Applying~\eqref{eq:ccc},

\[
\CCC = \frac{W}{r} = \frac{2{,}895{,}327}{28{,}150{,}350}\times 365 = 37.5\ \text{days}.
\]

\subsection{\texorpdfstring{$\CCC$}{CCC} by medium}

The same calculation by medium shows a difference of two orders of magnitude.

\begin{center}
\begin{tabular}{@{}lrrr@{}}
\toprule
Medium & $-\kappa_C$ (million yen) & $r$ (million yen) & $\CCC$ (days) \\
\midrule
Magnetic & 562{,}980 & 107{,}925 & 1{,}904 \\
Paper & 1{,}127{,}149 & 486{,}732 & 845 \\
Server & 670{,}552 & 14{,}331{,}607 & 17 \\
IC card & 534{,}645 & 13{,}224{,}087 & 15 \\
\bottomrule
\end{tabular}
\captionof{table}{$-\kappa_C$, $r$ and $\CCC$ by medium.}
\label{tab:media-ccc}
\end{center}

\subsection{\texorpdfstring{$\kappa$}{kappa} in the corporate statistics}

Applying~\eqref{eq:kappa} to the corporate sector as a whole, FY2024, all industries
excluding finance and insurance.

\begin{align}
\kappa_C &= \mathrm{DSO} = 71\ \text{days (capital 1bn yen and over)},\quad 35\ \text{days (under 10m yen)}\\
\kappa_S &= -\mathrm{DPO} = -44\ \text{days (1bn and over)},\quad -21\ \text{days (under 10m)}\\
\sum_i \kappa_i &= 27\ \text{days (1bn and over)},\quad 14\ \text{days (under 10m)}
\end{align}

\textbf{The sign of $\kappa$ depends on both size and industry, and the industry effect
dominates.} In retail, $\sum_i\kappa_i = -16$ days for firms of 1bn yen and over: the sign
reverses.

\section{Measuring the growth constraint}
\label{sec:gstar-measured}

The $g^\star=(m+d-I/r)/\CCC$ of Corollary~\ref{cor:gstar} needs four quantities. The
measurement proceeds in two stages.

\begin{center}
\begin{tabularx}{\textwidth}{@{}lXl@{}}
\toprule
Object & Quantities used & Period \\
\midrule
Long-run trend & $m$ and $\CCC$ only; the ratio $m/\CCC$ & FY2000--FY2024 \\
Level & $m$, $d$, $I$, $\CCC$; $g^\star$ itself & FY2016--FY2024 \\
\bottomrule
\end{tabularx}
\captionof{table}{The two stages of measuring the growth constraint: object and period.}
\label{tab:gstar-measurement-stages}
\end{center}

$d$ and $I$ are not used on the long-run side because depreciation and fixed-asset balances
were not obtained back to the 1990s. For $m$, the operating margin is used: the $m$ of
Chapter~\ref{sec:growth} is what funds increases in working capital, so a figure net of
selling and administrative expenses is closer to the definition.

\subsection{The corporate sector over time ($m/\CCC$)}

\begin{center}
\begin{tabular}{@{}rrrr@{}}
\toprule
Fiscal year & $\CCC$ (days) & $m$ & $m/\CCC$ \\
\midrule
2000 & 24.9 & 2.62\% & 38.5\% \\
2005 & 24.4 & 3.16\% & 47.3\% \\
2009 & 28.0 & 2.01\% & 26.2\% \\
2015 & 30.6 & 3.95\% & 47.1\% \\
2020 & 35.4 & 3.06\% & 31.5\% \\
2024 & 37.2 & 5.01\% & 49.2\% \\
\bottomrule
\end{tabular}
\captionof{table}{$\CCC$, $m$ and $m/\CCC$ for the corporate sector over time (all industries).}
\label{tab:mccc-trend}
\end{center}
\begin{center}\footnotesize All industries excluding finance and insurance, all sizes. $d$ and $I$ not included\end{center}

The trend is shown in Figure~\ref{fig:gstar-trend}.

\begin{figure}[htbp]
\centering
\insertfig{gstar-trend}
\caption{$m/\CCC$ over time (all industries, on an operating-margin basis). It differs from the $g^\star$ of Corollary~\ref{cor:gstar} by $d$ and $I$.}
\label{fig:gstar-trend}
\end{figure}

\begin{proposition}[$m/\CCC$ has not declined]
\label{prop:gstar-flat}
Between FY2000 and FY2024, $\CCC$ worsened by 49\%, from 24.9 to 37.2 days. But $m$ improved
by 91\%, from 2.62\% to 5.01\%, so
\[
\frac{m}{\CCC} : \quad 38.5\% \;\longrightarrow\; 49.2\%
\]
and the ratio in fact rose. The mean over 25 years is 41.1\% with a standard deviation of 7.3
points. Against a mean of 38.0\% over the first twelve years, the last thirteen average
44.0\%. No declining trend is present.
\end{proposition}

\begin{remark}[This proposition is not about $g^\star$ itself]
\label{rem:gstar-flat-caveat}
The $g^\star$ of Corollary~\ref{cor:gstar} is $(m+d-I/r)/\CCC$.
Proposition~\ref{prop:gstar-flat} concerns only $m/\CCC$, which excludes $d$ and $I$.

The reason for using $m/\CCC$ in the long time series is given in
Table~\ref{tab:gstar-measurement-stages}. The $g^\star$ including $d$ and $I$ is computed for
FY2016--FY2024 in Section~\ref{sec:gstar-measured}.

The proposition is therefore \textbf{a claim about the trend of a ratio}, not about the level
of the self-financeable growth rate. If $I$ has trended, the path of $g^\star$ can differ
from that of $m/\CCC$.
\end{remark}

\begin{remark}[Numerator and denominator must be measured separately]
Looking at the deterioration of $\CCC$ alone suggests a worsening growth constraint. But
Corollary~\ref{cor:gstar} is a ratio, and \textbf{no conclusion follows from a change in the
denominator alone}. Arguing from $\CCC$ without obtaining $m$ reaches the opposite of the
truth.

What drives the movement of $m/\CCC$ is mainly $m$. The troughs are FY2008--FY2009 (around
26\%) and FY2020 (31.5\%), both due to falling margins rather than to $\CCC$.
\end{remark}

\subsection{\texorpdfstring{$g^\star$}{g*} including depreciation and investment}

Corollary~\ref{cor:gstar} needs, besides $m$, the depreciation rate $d$ and investment $I$.
As Remark~\ref{rem:no-steady-investment} states, $I\approx d\,r$ is not assumed. $I$ is
estimated as the change in fixed-asset balances plus depreciation. Land is excluded because
it is not depreciated.
\begin{equation}
I \approx \bigl(\text{tangible}+\text{intangible}-\text{land}\bigr)_t
 - \bigl(\text{tangible}+\text{intangible}-\text{land}\bigr)_{t-1} + \text{depreciation}_t
\label{eq:I-estimate}
\end{equation}

\begin{center}
\begin{tabular}{@{}lrrrrrrr@{}}
\toprule
Size (capital) & $\CCC$ & $m$ & $d$ & $I/r$ & $g^\star$ (2024) & 9-yr mean & SD \\
\midrule
1bn yen and over & 47 & 7.8\% & 3.1\% & 4.2\% & 52.0\% & 49.3\% & 6.0 \\
100m -- 1bn & 35 & 4.8\% & 2.1\% & 4.3\% & 27.1\% & 39.6\% & 12.2 \\
50m -- 100m & 30 & 3.9\% & 1.8\% & 1.9\% & 45.8\% & 33.0\% & 15.2 \\
20m -- 50m & 25 & 3.5\% & 2.1\% & 1.2\% & 64.3\% & 35.0\% & 17.0 \\
10m -- 20m & 37 & 1.8\% & 2.3\% & 5.5\% & \textminus 13.8\% & 15.5\% & 30.8 \\
Under 10m & 26 & 1.5\% & 2.9\% & 2.0\% & 34.7\% & \textbf{3.1\%} & 41.4 \\
\bottomrule
\end{tabular}
\captionof{table}{$\CCC$, $m$, $d$, $I/r$ and $g^\star$ by capital class (FY2024, nine-year mean, SD).}
\label{tab:gstar-by-size}
\end{center}
\begin{center}\footnotesize All industries excluding finance and insurance; $\CCC$ in days; SD is the nine-year standard deviation of $g^\star$.
Ten years of data (FY2015--FY2024) were obtained, but because the $I$ of~\eqref{eq:I-estimate} needs the previous year's fixed-asset balance, $g^\star$ can be computed only from FY2016, and the mean and SD are taken over the nine years FY2016--FY2024\end{center}

\begin{remark}[Single-year values differ in reliability by size]
\label{rem:gstar-reliability}
The nine-year standard deviation of $I/r$ is 0.3 for firms of 1bn yen and over and 2.9 for
those under 10m --- an order of magnitude apart. Because~\eqref{eq:I-estimate} uses changes in
balances, retirements and disposals and \textbf{movement between capital classes} contaminate
it. Capital changes with equity issuance irrespective of performance, so firms move between
classes.

The direction of movement is not constant. The $I/r$ of the under-10m class recorded
$-1.6\%$ in 2016 and $+7.5\%$ in 2022, so this is not a one-directional bias. Averaging
therefore largely cancels the effect, but single-year values are unreliable for the small
classes.

This is why the table reports both single-year and nine-year figures with standard
deviations.
\end{remark}

Following Remark~\ref{rem:gstar-reliability}, judgments use the nine-year means.

\begin{proposition}[The constraint on small firms is not $\CCC$]
\label{prop:small-margin}
On nine-year means, $g^\star$ for firms under 10m yen of capital is 3.1\%, more than an order
of magnitude below the 49.3\% of firms of 1bn and over. The difference does not come from
$\CCC$: the $\CCC$ of the small classes (26--37 days) is in fact shorter than that of large
firms (47 days).

The cause is the numerator. $m+d-I/r$ is 6.7\% for firms of 1bn and over and 2.4\% for those
under 10m; besides a low $m$ of 1.5\%, an investment burden exceeding $d$ contributes.
\end{proposition}

In the single year, the under-10m class shows a high $34.7\%$, but with a standard deviation
of 41.4 no meaning should be read into that value.

\subsection{\texorpdfstring{$g^\star$}{g*} by industry}

\begin{center}
\begin{tabular}{@{}lrrrrrr@{}}
\toprule
Industry & $\CCC$ & $m$ & $d$ & $I/r$ & 9-yr mean $g^\star$ & SD \\
\midrule
Manufacturing & 42 & 5.4\% & 2.8\% & 4.2\% & 34.8\% & 5.6 \\
All industries & 37 & 5.0\% & 2.5\% & 3.5\% & 37.7\% & 7.5 \\
Wholesale & 33 & 2.2\% & 0.7\% & 0.4\% & 20.9\% & 6.9 \\
Construction & 40 & 4.5\% & 1.5\% & 1.4\% & 47.1\% & 11.3 \\
Retail & 23 & 2.7\% & 1.5\% & 1.6\% & 29.0\% & 12.2 \\
Employment placement and dispatch & 33 & 3.9\% & 0.6\% & \textminus 0.3\% & 42.7\% & 13.9 \\
Information and communications & 58 & 8.2\% & 4.8\% & 5.2\% & 61.2\% & 16.5 \\
Other services & 34 & 5.4\% & 3.0\% & 5.5\% & 47.9\% & 32.3 \\
\midrule
Accommodation & 22 & 6.6\% & 3.5\% & 6.3\% & \textminus 204.2\% & \textbf{872} \\
Food services & 5 & 1.2\% & 2.4\% & \textminus 0.4\% & 266.9\% & \textbf{3334} \\
\bottomrule
\end{tabular}
\captionof{table}{$\CCC$, $m$, $d$, $I/r$ and the nine-year mean $g^\star$ by industry, with standard deviations.}
\label{tab:gstar-by-industry}
\end{center}

The bottom two industries in Table~\ref{tab:gstar-by-industry} exhibit the divergence of
Remark~\ref{rem:gstar-divergence}. The $\CCC$ of food services is 5 days, so a one-point
movement in the numerator moves $g^\star$ by about 73 points. The standard deviation of 3334
is the consequence of that amplification:
\textbf{$g^\star$ is not functioning as an indicator}.

\begin{remark}[Two meanings of a short $\CCC$]
In the types of Part~\ref{part:catalog}, food services is family 1-1 (immediate exchange) and
$\tau_i\approx 0$ is a structural property (Corollary~\ref{cor:ccc-level}).
A short $\CCC$ indicates sound liquidity, but
\textbf{it simultaneously makes the ratio $g^\star$ meaningless}.
Corollary~\ref{cor:gstar} is useful only for objects whose $\CCC$ is some tens of days or
more.
\end{remark}

\subsection{The $m$ of small firms is decided by directors' remuneration}
\label{sec:officer-comp}

Proposition~\ref{prop:small-margin} showed that $m$ is the constraint but not why $m$ is low.
Decompose the cost structure.

\begin{center}
\begin{tabular}{@{}lrrrr@{}}
\toprule
Size (capital) & Gross margin & SG\&A ratio & Personnel ratio & $m$ \\
\midrule
1bn yen and over & 23.6\% & 15.8\% & 7.7\% & 7.8\% \\
100m -- 1bn & 21.5\% & 16.8\% & 10.3\% & 4.8\% \\
50m -- 100m & 25.9\% & 22.0\% & 12.3\% & 3.9\% \\
20m -- 50m & 26.0\% & 22.4\% & 15.0\% & 3.5\% \\
10m -- 20m & 33.8\% & 32.0\% & 17.7\% & 1.8\% \\
Under 10m & \textbf{46.3\%} & \textbf{44.8\%} & 22.7\% & 1.5\% \\
\bottomrule
\end{tabular}
\captionof{table}{Gross margin, SG\&A ratio, personnel ratio and $m$ by capital class (FY2024).}
\label{tab:cost-structure-by-size}
\end{center}
\begin{center}\footnotesize All industries excluding finance and insurance, FY2024\end{center}

\textbf{Small firms have high gross margins}: 46.3\% against 23.6\% for large firms. Their
SG\&A ratio, however, is 44.8\% against 15.8\%. As Remark~\ref{rem:cadmin-negligible} notes,
part of that difference may be $C_{\mathrm{admin}}$: at a scale where invoicing and chasing
are done by hand, the operating cost of contracts is booked in SG\&A. This text does not
separate it. The low $m$ is therefore not due to a lack of pricing power; the drop occurs at
the SG\&A stage.

Moving personnel costs back into cost of sales narrows the gap in gross margin from 22.7
points to 7.7 (15.9\% against 23.6\%). Two-thirds of the gap was a difference of accounting
classification. Even after the adjustment, however, the small classes remain higher.

Decomposing personnel costs into directors and employees identifies the cause.

\begin{center}
\begin{tabular}{@{}lrrrr@{}}
\toprule
Size (capital) & Directors' remuneration & Employee wages & Directors' share & $m$ \\
\midrule
1bn yen and over & 0.16\% & 7.6\% & 2.0\% & 7.8\% \\
100m -- 1bn & 0.40\% & 9.9\% & 3.9\% & 4.8\% \\
50m -- 100m & 0.99\% & 11.3\% & 8.1\% & 3.9\% \\
20m -- 50m & 2.35\% & 12.7\% & 15.6\% & 3.5\% \\
10m -- 20m & 3.98\% & 13.7\% & 22.5\% & 1.8\% \\
Under 10m & \textbf{8.04\%} & 14.6\% & \textbf{35.5\%} & 1.5\% \\
\bottomrule
\end{tabular}
\captionof{table}{Directors' remuneration, employee wages, directors' share and $m$ by capital class.}
\label{tab:officer-comp-by-size}
\end{center}
\begin{center}\footnotesize The directors' share is directors' remuneration as a fraction of personnel costs\end{center}

\begin{proposition}[Directors' remuneration explains the difference in $m$ by size]
\label{prop:officer}
Directors' remuneration as a fraction of sales is 0.16\% for firms with capital of 1bn yen
and over and 8.04\% for those under 10m --- a \textbf{fiftyfold difference}. The gap in $m$
between the largest and smallest classes is 6.3 points, while the gap in directors'
remuneration is 7.88 points, large enough to account for it.

Adding directors' remuneration back into operating profit removes the monotone ordering.
\begin{center}
\begin{tabular}{@{}lrr@{}}
\toprule
Size & $m$ & After adding remuneration back \\
\midrule
1bn yen and over & 7.8\% & 8.0\% \\
20m -- 50m & 3.5\% & 5.9\% \\
Under 10m & 1.5\% & \textbf{9.5\%} \\
\bottomrule
\end{tabular}
\captionof{table}{$m$ before and after adding directors' remuneration back into operating profit, by size.}
\label{tab:officer-comp-adjusted}
\end{center}
The smallest class becomes the highest. The same holds within industries: $m$ is higher for
the small classes in 1 of 9 industries as reported, but in 7 after remuneration is added
back.
\end{proposition}

\begin{remark}[A discrepancy with the definition of $m$]
Chapter~\ref{sec:growth} defined $m$ as what funds increases in working capital. If
directors' remuneration is discretionary, what a small operator can actually use is not
operating profit but \textbf{operating profit plus directors' remuneration}. On that basis
the numerator for the under-10m class is not the 2.4\% of $m+d-I/r$ but 10.4\% once
remuneration of 8.04\% is added. Dividing by $\CCC=26$ days gives 146\%, far above the 51.7\%
of large firms (Section~\ref{sec:gstar-measured}; nine-year mean 49.3\%).

That is, Proposition~\ref{prop:small-margin}'s claim that small firms are constrained by $m$
is correct only so long as reported operating profit is taken to be $m$;
\textbf{with an $m$ faithful to the definition, the conclusion reverses}.
\end{remark}

\subsection{Comparison with unincorporated businesses}
\label{sec:individual-firms}

Proposition~\ref{prop:officer} rests on the corporate statistics. Whether the same structure
appears in a different statistic, with different size classes, and in a sample including
\textbf{entities without corporate personality}, is checked here.

The Basic Survey on Small and Medium Enterprises was used: a sample of about 110{,}000 firms
in which four corporate size classes by headcount and unincorporated businesses appear in the
same table. Unincorporated businesses have no line item for directors' remuneration, so the
proprietor's return to labour remains in operating profit.

\begin{quote}
\textbf{Registered implication}: the operating margin of unincorporated businesses is close
to the corporate small classes' figure ``after adding remuneration back''. The two are
substantively the same economic agent and differ only in accounting classification.
\end{quote}

\begin{center}
\begin{tabular}{@{}lrrrrrrrr@{}}
\toprule
Class & 2019 & 2020 & 2021 & 2022 & 2023 & 2024 & Mean & SD \\
\midrule
Corporate, 5 or fewer & 2.7 & 1.1 & 2.6 & 3.0 & 2.6 & 3.0 & 2.5 & 0.7 \\
Corporate, 6--20 & 2.4 & 1.5 & 2.5 & 2.8 & 3.0 & 3.2 & 2.6 & 0.6 \\
Corporate, 21--50 & 2.8 & 2.1 & 3.2 & 2.8 & 3.4 & 3.7 & 3.0 & 0.6 \\
Corporate, 51 and over & 3.4 & 2.8 & 3.2 & 3.6 & 3.6 & 4.0 & 3.4 & 0.4 \\
Unincorporated & 17.0 & 18.5 & 18.2 & 15.8 & 15.4 & 16.3 & \textbf{16.9} & 1.3 \\
\bottomrule
\end{tabular}
\captionof{table}{Operating margin by corporate size class and for unincorporated businesses (FY2019--FY2024).}
\label{tab:individual-vs-corp-margin}
\end{center}
\begin{center}\footnotesize Operating margin (\%), all industries, by fiscal year\end{center}

\textbf{The implication is supported.} The operating margin of unincorporated businesses is
about five times the corporate figure in all six years, with a standard deviation of 1.3.
This is not a single-year peculiarity. The corporate figure after adding remuneration back in
Section~\ref{sec:officer-comp} was 9.5\%. The 16.9\% of unincorporated businesses exceeds
that, but is far closer to it than to the 1.5--3.0\% before adding back.

The personnel-cost ratio explains the difference.

\begin{center}
\begin{tabular}{@{}lrrr@{}}
\toprule
Class & Operating margin & Personnel ratio & Total \\
\midrule
Corporate, 5 or fewer & 2.5\% & 14.7\% & 17.2\% \\
Corporate, 6--20 & 2.6\% & 16.9\% & 19.5\% \\
Unincorporated & 16.9\% & \textbf{9.8\%} & 26.7\% \\
\bottomrule
\end{tabular}
\captionof{table}{Operating margin and personnel ratio, corporate and unincorporated (six-year means).}
\label{tab:individual-vs-corp-laborcost}
\end{center}
\begin{center}\footnotesize Six-year means. The personnel ratio is labour cost in cost of sales plus personnel cost in SG\&A\end{center}

\textbf{Unincorporated businesses have the lowest personnel ratio}, because the proprietor's
labour is not booked as a cost. What appears in personnel costs as directors' remuneration in
a company remains in operating profit for an unincorporated business.

\begin{remark}[The sign reverses during the pandemic]
\label{rem:covid-reversal}
In FY2020 the operating margin fell in every corporate class (2.7\% $\to$ 1.1\% for firms of
five or fewer). \textbf{Only unincorporated businesses rose} (17.0\% $\to$ 18.5\%).

This has the same form as the observation in Remark~\ref{rem:officer-undecidable}. Because
the operating profit of an unincorporated business contains the proprietor's own take, the
ratio does not fall when sales fall unless the proprietor reduces that take. If employee
wages and outsourcing are cut first, it rises.

That directors' remuneration in companies moves independently of profit, and that the
operating margin of unincorporated businesses rises in a downturn, are
\textbf{two appearances of one and the same phenomenon}.
\end{remark}

\begin{remark}[Capital expenditure by unincorporated businesses]
\label{rem:individual-capex}
The same survey gives $I/r$ for unincorporated businesses: a six-year mean of 1.9\% with a
standard deviation of 1.0. That is lower and less stable than the 3.1\% (SD 0.1) of corporate
firms with 51 or more employees. The share of firms making any capital expenditure is 8.1\%
(FY2024): more than nine in ten invest nothing.

Chapter~\ref{ch:solo} assumed $A\approx 0$ for a solo business. The assumption agrees with
the data on unincorporated businesses.
\end{remark}

\begin{remark}[What cannot be decided]
\label{rem:officer-undecidable}
Whether directors' remuneration is a discretionary appropriation of profit or a return to
labour cannot be decided here.

An observation supporting discretion: the directors' remuneration ratio of the under-10m
class stayed in the narrow band 7.6--9.1\% throughout FY2015--FY2024, while $m$ over the same
period ranged from $-2.0\%$ to $+2.1\%$.
\textbf{Profit moves and remuneration does not.}

But the labour-return account cannot be rejected either. Directors of small firms work in
substance as employees, and it is equally possible that the remuneration is the return to
that. Distinguishing the two needs remuneration per director and hours worked, but
\textbf{directors are not employees under the Labour Standards Act and no statistics on their
hours exist}. Both the Monthly Labour Survey and the Basic Survey on Wage Structure are
confined to employees. This belongs to category (i), no measurement exists, of
Section~\ref{sec:obstacle-taxonomy}.
\end{remark}

\subsection{By industry (FY2024)}

\begin{center}
\begin{tabular}{@{}lrrr@{}}
\toprule
Industry & $\CCC$ (days) & $m$ & $g^\star$ \\
\midrule
Pure holding companies & 12 & 57.7\% & 1710\% \\
Water transport & 7 & 4.8\% & 261\% \\
Mining, quarrying and gravel extraction & 29 & 19.0\% & 236\% \\
Motor vehicles and parts & 9 & 5.1\% & 201\% \\
\midrule
Leasing & 210 & 6.6\% & 11.4\% \\
Petroleum and coal products & 20 & 0.6\% & 11.4\% \\
Textiles & 78 & 2.0\% & 9.2\% \\
Fisheries & 40 & \textminus 1.5\% & negative \\
Agriculture and forestry & 22 & \textminus 4.0\% & negative \\
\bottomrule
\end{tabular}
\captionof{table}{$\CCC$, $m$ and $g^\star$ by industry (FY2024; the extremes).}
\label{tab:gstar-industry-extremes}
\end{center}

Where the numerator $m+d-I/r$ is negative, $g^\star<0$, and by
Proposition~\ref{prop:fcf} \textbf{free cash flow is negative even at zero growth}.

A long $\CCC$ and a small $g^\star$ do not coincide
(Table~\ref{tab:gstar-industry-extremes}). Leasing has the longest $\CCC$ at 210 days but a
high $m=6.6\%$, so on $m$ alone it sits at the same level as petroleum and coal products
($\CCC=20$ days, $m=0.6\%$).
\textbf{$\CCC$ alone cannot decide the growth constraint.}
The $g^\star$ including $d$ and $I$ is given in Section~\ref{sec:gstar-measured}.

\section{Right and exercise, and cognitive surplus}

The $\phi_{\rm cog}=p\cdot\E[\dbar-\delta]$ of~\eqref{eq:cog} is a measurable quantity, and
measurements exist. Four are given below.
\textbf{What this text obtained itself, however, is only an upper bound; the value is
unknown.}

\subsection{\texorpdfstring{$\phi_{\rm cog}$}{phi-cog} is measurable}
\label{sec:phicog-measurable}

\paragraph{This text's measurement: family 2 recoveries on expiry.}

\begin{align}
p\cdot\E[\dbar-\delta] &\leq 20{,}634\ \text{million yen}\\
\frac{20{,}634}{2{,}895{,}327} &= 0.71\%\quad\text{(of unused balances)}\\
\frac{20{,}634}{28{,}150{,}350} &= 0.073\%\quad\text{(of issuance)}
\end{align}

The inequality holds because this figure also contains recoveries through refunds.
\textbf{What was obtained is an upper bound, not $\phi_{\rm cog}$ itself.}
The three examples below all give the value itself.

\paragraph{Gym memberships.}

\cite{dellavigna2006} matched the contract choices and daily attendance records of 7{,}752
members across three US gyms, measuring the $\dbar-\delta$ of~\eqref{eq:cog} directly.
Members predicted 9.50 visits a month and made 4.17. Members on plans above \$70 a month
attended on average only 4.3 times, paying more than \$17 a visit where a ten-visit pass
would cost \$10. The loss over a membership averaged \$600 against total payments of about
\$1{,}400, so $\phi_{\rm cog}$ is about 43\% of the amount paid.

\paragraph{Subscription renewals.}

\cite{einav2025} covers a wider range. Using comprehensive data from a payment-card network,
it exploits as an exogenous shock the fact that card renewal forces an active renewal
decision. Across ten subscriptions, \textbf{inattention raises sellers' revenue by 89\% on
average} (estimates ranging from 14\% to over 200\%). The attention parameter $\lambda$ is
estimated at 0.18 on average, ranging from 0.04 to 0.50 across services.

\paragraph{Breakage on prepaid balances.}

\cite{nber2026pnbl} treats family 2-4 (prepaid) directly. Analysing transaction data on over
four million people for a scheme adding a bonus to prepaid restaurant balances, it records
that \textbf{about 40\% of prepaid value goes unused}. Operators earn a median of \$5.50 of
breakage profit per dollar of bonus, and the profit from breakage exceeds the revenue from
increased spending.

\begin{proposition}[Three conditions for measuring $\phi_{\rm cog}$]
\label{prop:phicog-requirements}
Estimating $\phi_{\rm cog}$ requires all three of the following simultaneously.
\begin{enumerate}[label=(\arabic*)]
\item \textbf{The contractual right $\dbar$}: the tariff and any usage ceiling
\item \textbf{Realized use $\delta$}: individual records for the same party; aggregates will not do
\item \textbf{An exogenous event forcing attention}: card renewal, a change of terms, a price revision
\end{enumerate}
\end{proposition}

\begin{proof}[Grounds]
Without (1) and (2) the difference $\dbar-\delta$ cannot be defined. Without (3), observed
continuation cannot be distinguished between ``continued because it was worth it'' and
``continued through inattention'', and $\lambda$ is not identified.
\cite{dellavigna2006} satisfies (1) and (2) through contract data and attendance records;
\cite{einav2025} satisfies (3) through card renewal.
\end{proof}

This text's family 2 measurement stopped at an upper bound \textbf{because (2) and (3) were
missing}. Aggregate statistics contain no individual records and there is no exogenous shock.
In the taxonomy of Section~\ref{sec:obstacle-taxonomy} this is not (i), no measurement exists,
but (iv), access constraint: the method of measurement is established and what is needed is
reaching the data.

\subsection{A detour through the menu}
\label{sec:phicog-detour}

Condition (2) of Proposition~\ref{prop:phicog-requirements}, individual realized use, is not
available in published statistics. But \textbf{a substitution of definitions can sometimes
provide a detour}.

When a flat rate $F$ and a per-use price $p$ coexist on the same menu, the break-even number
of uses can be computed.
\begin{equation}
\delta^{*} = F / p
\label{eq:delta-star}
\end{equation}
Even where $\dbar$ is ``unlimited'' and unmeasurable, \eqref{eq:delta-star} follows
\textbf{from the menu alone}. Using it as a proxy for $\dbar$, all that is needed is the
industry-average number of uses $\delta$.

Since~\eqref{eq:delta-star} is equivalent to $p\,\delta^{*}=F$, substituting
into~\eqref{eq:cog} gives
\begin{equation}
\phi_{\mathrm{cog}} \approx p\,\E[\delta^{*}-\delta] = F - p\,\E[\delta]
\label{eq:cog-menu}
\end{equation}
\textbf{Under this proxy, cognitive surplus is the difference between the flat fee and the
per-use equivalent.} The registered implication $\delta^{*}-\delta>0$ becomes the same claim
as \eqref{eq:cog-menu} being positive, that is, the flat fee exceeding the per-use equivalent
of average use.

\begin{quote}
\textbf{Registered implication}: for flat-rate members, $\delta^{*} - \delta > 0$; the
average member does not reach break-even.
\end{quote}

\paragraph{Result.}
From the Survey on Selected Service Industries, $\delta$ is 89 visits a year per individual
member, that is 7.42 a month. $\delta^{*}$ was computed from the published prices of one
large full-service operator (2026, tax included).

\begin{center}
\begin{tabular}{@{}lrrrr@{}}
\toprule
Facility category & Unlimited plan & Per visit & $\delta^{*}$ & $\delta^{*}-\delta$ \\
\midrule
I & 12{,}680 & 2{,}530 & 5.01 & \textminus 2.40 \\
II & 13{,}980 & 2{,}860 & 4.89 & \textminus 2.53 \\
III & 16{,}580 & 3{,}190 & 5.20 & \textminus 2.22 \\
IV & 19{,}080 & 3{,}630 & 5.26 & \textminus 2.16 \\
\bottomrule
\end{tabular}
\captionof{table}{Unlimited and per-visit prices by facility category, and the break-even number of visits $\delta^{*}$.}
\label{tab:predk-breakeven}
\end{center}

\textbf{The implication is rejected.} The sign is negative in all four categories
(Table~\ref{tab:predk-breakeven}): for the average member the flat contract is rational.

The effective unit price --- total dues divided by total visits --- points the same way,
falling 29\% over 25 years from 1{,}147 yen in 2000 to 815 yen in 2023. Against a per-visit
market rate of 2{,}530--3{,}630 yen, \textbf{members are using the facility more cheaply than
paying per visit}.

\begin{remark}[When the detour applies]
\label{rem:detour-limit}
Equation~\eqref{eq:delta-star} is uniquely determined only when the menu is a binary choice
between flat and per-use. Actual tariffs are stepped: the operator above offers four visits a
month, eight a month, and unlimited. The threshold then changes with the comparison chosen.
Choosing the unlimited plan is rational at 5.0 visits a month against per-visit pricing, but
requires 9.1--9.5 against the eight-visit plan.
\textbf{Average use of 7.42 is rational on the former and irrational on the latter.}

The detour therefore applies only to binary menus; with a stepped menu the proxy for $\dbar$
is not uniquely determined.
\end{remark}

\begin{remark}[$\phi_{\rm cog}=0$ does not follow from this section]
The rejection is a result about the mean. The findings of \cite{dellavigna2006} concern the
distribution, and a structure in which the lower tail loses is not excluded by a mean above
break-even. Mean and distribution are different claims.
\end{remark}

\begin{remark}[What remains unmeasured within this framework]
The prior work concentrates on subscriptions and memberships, that is, family 2-1 of
Part~\ref{part:catalog}. No comparable estimates were found for family 2-4 (prepaid, gift
cards) or family 2-5 (options and warranties). Nor was any study found estimating the
\emph{share} of $\phi_{\rm cog}$ within total $\phi$ across industries. What
\cite{einav2025} measures is the increment to revenue, not a composition of profit.
\end{remark}

The reason for purchase by medium serves as a proxy for the transaction frequency $\nu$.

\begin{center}
\begin{tabular}{@{}lrrr@{}}
\toprule
Medium & For own use & As a gift & $\CCC$ (days) \\
\midrule
IC card & 95.5\% & 0.0\% & 15 \\
Server (online) & 78.1\% & 14.0\% & 17 \\
Paper & 41.5\% & 30.1\% & 845 \\
\bottomrule
\end{tabular}
\captionof{table}{Reason for purchase by medium (own use versus gift) and $\CCC$.}
\label{tab:media-purchase-reason}
\end{center}

The higher the gift share, the lower $\nu$ and the longer $\CCC$. The $\lambda(\nu)$
of~\eqref{eq:decay} itself was not measured.

\section{The information structure \texorpdfstring{$\mathcal{G}$}{G} made concrete}

The $\mathcal{G}$ of Chapter~\ref{sec:info} --- verifiable shared information --- is
observable in family 5 as a concrete set of fields. The items published on the Employment
Service Information Site are what $\mathcal{G}$ consists of.

\begin{center}
\begin{tabularx}{\textwidth}{@{}lXl@{}}
\toprule
Symbol & Instance & Status \\
\midrule
$y$ & separations within six months among those placed on open-ended contracts & numeric, six years \\
Breadth of $\mathcal{G}$ & the count of ``separation not ascertained'' (the smaller, the wider $\mathcal{G}$) & numeric; zero in all years for both observations \\
$b$ (outcome-contingency) & refund schemes & presence is binary; rates and periods are not in a uniform format \\
Proxy on the price side & fee rates achieved by occupation & continuous; 307 firms (Section~\ref{sec:predX-result}) \\
$\sigma^2$ (noise in the outcome) & baseline separation rates by occupation & \textbf{not measured} \\
\bottomrule
\end{tabularx}
\captionof{table}{The information structure $\mathcal{G}$ made concrete, and the observation status of each symbol (family 5).}
\label{tab:info-structure-observed}
\end{center}

Measured values of $y$ at the two establishments that could be confirmed:

\[
\frac{y}{\text{open-ended placements}} = \frac{124}{4{,}348} = 2.85\%, \qquad
\frac{3}{189} = 1.59\% .
\]

Testing the $b^\star=1/(1+\gamma\sigma^2 k)$ of Example~\ref{ex:bstar} requires both $b$ and
$\sigma^2$ as continuous quantities, but refund rates and periods are not in a uniform format
and the uncontrollable part cannot be extracted from $\sigma^2$.

A test on the price side was possible, however. Fee rates achieved by occupation were
obtained as a continuous quantity for 307 firms and correlated with separation rates
(Section~\ref{sec:predX-result}).

\section{The three-way decomposition of surplus}

The three components of~\eqref{eq:surplus} can be separated only in the case of
Chapter~\ref{sec:platform-fee}.

\begin{center}
\begin{tabularx}{\textwidth}{@{}lrX@{}}
\toprule
Component & Measured & Grounds \\
\midrule
$\phi_{\rm prod}$ (price of function) & 20.0 points & within-platform comparison on Gumroad, differing only in whether customer acquisition is included \\
$\phi_{\rm barg}$ (rent from market power) & 15.0 points & the App Store's two size tiers; function identical \\
$\phi_{\rm risk}$ (price of transfer and pooling) & --- & $\approx 0$ for these two firms; positive where cases are screened individually \\
$\phi_{\rm cog}$ & --- & not measured in this domain \\
\bottomrule
\end{tabularx}
\captionof{table}{Measured values of the three-way decomposition of surplus, and their grounds (platform fees).}
\label{tab:surplus-decomposition-measured}
\end{center}

This is \textbf{the only case in which two of the three components could be separated
numerically for a single transaction}. No procedure for separating the components exists in
the other domains.

The two firms are intermediaries operating at uniform rates published in a price list, so the
risk component $\phi_{\rm risk}$ was negligible and separation was possible. Where an
intermediary sets rates by individual screening, a third term is added and the separation
fails (Remark~\ref{rem:fee-incomplete}).

\section{The terms of the launch cost}

For each term of~\eqref{eq:E-required-2}, the price of externalization was measured, as the
difference in effective rate on a \$50 transaction.

\begin{center}
\begin{tabular}{@{}llr@{}}
\toprule
Term & Means of externalization & Price \\
\midrule
$T_{\rm inst}$ & payment processing & 3.6\% \\
$T_{\rm inst}$ (tax) & merchant of record & +1.4--2.4 points \\
$T_{\rm dev}$ (delivery) & a sales platform with hosting & +8--9 points \\
$T_{\rm acq}$ & marketplace customer acquisition & +20 points \\
$T_{\rm credit}$ & self-funded through a free period & the cost of $\kappa>0$ \\
\bottomrule
\end{tabular}
\captionof{table}{The terms of the launch cost, the means of externalization, and the price.}
\label{tab:launch-cost-externalization}
\end{center}

\textbf{$T_{\rm acq}$ exceeds the sum of all the other items.}
$T_{\rm dev}$ itself (the development period) and the enumeration of $\hat\Omega$ cannot be
externalized and therefore have no price.

\section{Quantities relating to selection and identification}

\begin{center}
\begin{tabularx}{\textwidth}{@{}llX@{}}
\toprule
Symbol & Measured & Source \\
\midrule
$\Prob(S=1)$ (listed) & $3{,}975 / 3{,}600{,}000 = 0.11\%$ & Section~\ref{sec:frame} \\
$\Prob(S=1)$ (mentioned) & of order $10^{-4}$ & as above \\
$\tau$ (insolvency time) & \textbf{not observed} & outside the sample through conditioning on survival \\
$(a,p,c)$ & \textbf{not identified} & Proposition~\ref{prop:apc} \\
Upper bound on $\lambda_{\rm novel}$ & 48{,}185 a year (published CVEs) & Chapter~\ref{part:unmanned} \\
Effective $\lambda_{\rm novel}$ & 245 a year (confirmed exploitation) $=0.51\%$ & as above \\
\bottomrule
\end{tabularx}
\captionof{table}{Measured quantities relating to selection and identification, and their sources.}
\label{tab:selection-identification-measures}
\end{center}

\section{Quantities that could not be measured}

The following are defined by the framework but could not be connected to real data
(Table~\ref{tab:unmeasured-quantities}).

\begin{center}
\begin{tabularx}{\textwidth}{@{}llX@{}}
\toprule
Symbol & Name & Why it cannot be connected \\
\midrule
$M(t)$ & cash balance & needs per-firm data; meaningless in aggregate \\
$\Omega,\ \mathcal{F}_t$ & state space, filtration & not observable in principle \\
$v$ & valuation map & subjective; the definition cannot be fixed \\
$\beta_i$ & discount factor by party & requires experimental methods \\
$\mathrm{Cap}$ & capacity ceiling & differs by operator and is not disclosed \\
$\sigma^2$ & noise in the outcome & variance by occupation is not tabulated (Section~\ref{sec:what-to-measure}) \\
$\lambda(\nu)$ & decay rate of cognitive surplus & the numerator $\phi_{\rm cog}$ is itself unmeasured \\
$\phi_{\rm cog}$ & cognitive surplus & only an upper bound here; the method is established (Section~\ref{sec:phicog-measurable}) \\
$\hat\Omega$ & envisaged state space & exists only inside the designer \\
$\tau$ & insolvency time & outside the sample through conditioning on survival (Proposition~\ref{prop:pathwise}) \\
& & \\
\bottomrule
\end{tabularx}
\captionof{table}{Quantities defined by the framework that could not be connected to real data.}
\label{tab:unmeasured-quantities}
\end{center}

$\phi_{\rm cog}$ is measurable in principle (Proposition~\ref{prop:phicog-requirements}); that
this text could not reach it is a matter of data granularity. The rest belong to (i), no
measurement exists, or are unobservable in principle.

\begin{remark}[The framework is larger than what can be measured]
About half the quantities defined could not be connected to real data. That is to say the
framework describes the world more finely than the world is recorded, and it shows that
precision of description and testability are independent. In the taxonomy of
Section~\ref{sec:obstacle-taxonomy}, most of what cannot be connected belongs to (i), no
measurement exists, and effort does not resolve it.
\end{remark}
